\documentclass[11pt]{article}

\usepackage[margin=1in]{geometry}
\usepackage{amsmath,amssymb}
\usepackage{booktabs,longtable}
\usepackage{hyperref}

\hypersetup{
  colorlinks=true,
  linkcolor=blue,
  urlcolor=blue
}

\title{Surface Pond Environment Model}
\author{LIFEORIG working model specification}
\date{\today}

\begin{document}
\maketitle

\section{Purpose}

The surface pond environment is a shallow aqueous reactor exposed to atmosphere,
UV radiation, sediment or mineral surfaces, and changing water depth. It is
different from a hydrothermal vent and different from volcanic rock pores.

The key physical idea is:
\begin{equation}
\boxed{
\text{shallow pond}
\rightarrow
\text{water-depth cycles}
\rightarrow
\text{concentration/dilution}
\rightarrow
\text{UV and sediment chemistry}
}
\end{equation}

The first implementation should be a single well-mixed pond. Spatial structure
can be added later.

\section{Geometry}

The first surface pond model is zero-dimensional: one well-mixed shallow water
body above a sediment or mineral surface.

\begin{verbatim}
atmosphere / UV input
        |
        v
+--------------------------+
| shallow pond water       |
| dissolved molecules      |
+--------------------------+
| sediment / mineral floor |
+--------------------------+
\end{verbatim}

The geometry is defined by:
\begin{equation}
V(t)=A_{\rm pond}h(t),
\end{equation}
where $V(t)$ is water volume, $A_{\rm pond}$ is pond area, and $h(t)$ is water
depth.

\section{State variables}

\begin{longtable}{lll}
\toprule
Symbol & Code name & Meaning \\
\midrule
$t$ & \texttt{time} & simulation time \\
$A_{\rm pond}$ & \texttt{area} & horizontal pond area \\
$h(t)$ & \texttt{water\_depth} & instantaneous water depth \\
$V(t)$ & \texttt{volume} & water volume, $A_{\rm pond}h(t)$ \\
$T$ & \texttt{temperature} & pond water temperature \\
$P$ & \texttt{pressure} & local pressure \\
$\mathrm{pH}$ & \texttt{pH} & acidity/basicity of pond water \\
$I$ & \texttt{ionic\_strength} & dissolved salt strength \\
$F_{\rm UV}$ & \texttt{uv\_flux} & UV forcing level \\
$A_{\rm sed}$ & \texttt{sediment\_surface\_area} & reactive sediment/mineral area \\
$C_i$ & \texttt{concentrations[species]} & concentration of species $i$ \\
\bottomrule
\end{longtable}

\section{Water-depth cycle}

The simplest water-depth cycle is periodic:
\begin{equation}
h(t)
=
h_0
+
\Delta h
\cos\left(\frac{2\pi t}{\tau_{\rm cyc}}\right),
\end{equation}
with clipping:
\begin{equation}
h(t)\ge h_{\rm min}.
\end{equation}

The variables mean:
\begin{itemize}
  \item $h_0$ is the mean water depth.
  \item $\Delta h$ is the amplitude of drying/refilling.
  \item $\tau_{\rm cyc}$ is the cycle period.
  \item $h_{\rm min}$ prevents the pond from becoming mathematically negative.
\end{itemize}

\section{Concentration and dilution}

When the pond dries, the same number of molecules occupies less water. The
concentration rises. When the pond refills, concentration falls.

If $h_{\rm old}$ is the previous depth and $h_{\rm new}$ is the new depth:
\begin{equation}
C_i^{\rm new}
=
C_i^{\rm old}
\frac{h_{\rm old}}{h_{\rm new}}.
\end{equation}

This is the central surface-pond mechanism. It is why shallow ponds can drive
prebiotic chemistry: second-order and higher-order reactions speed up strongly
when concentrations increase during drying.

\section{Chemical update}

The minimal concentration equation is:
\begin{equation}
\frac{dC_i}{dt}
=
R_i(\mathbf{C},T,\mathrm{pH},I,F_{\rm UV},A_{\rm sed})
+
S_i^{\rm atm}
-
k_{i}^{\rm photo}F_{\rm UV}C_i
-
k_{\rm wash}C_i
-
\frac{C_i}{h}\frac{dh}{dt}.
\end{equation}

The terms mean:
\begin{itemize}
  \item $R_i$ is production or loss from the chemical reaction network.
  \item $S_i^{\rm atm}$ is atmospheric input or deposition of species $i$.
  \item $k_i^{\rm photo}F_{\rm UV}C_i$ is UV photolysis loss.
  \item $k_{\rm wash}C_i$ is washout or runoff loss.
  \item $-(C_i/h)(dh/dt)$ is concentration during drying or dilution during
  refilling.
\end{itemize}

\section{Code structure}

The first implementation is:
\begin{verbatim}
SurfacePond
    state: PondState
    concentrations
    atmospheric_input
    photolysis_rates
    water_depth(t)

PondState
    time
    water_depth
    temperature
    pressure
    pH
    ionic_strength
    uv_flux
    sediment_surface_area
    volume
    concentrations
\end{verbatim}

\section{Comparison with other environments}

\begin{longtable}{lll}
\toprule
Environment & Main driver & Main geometry \\
\midrule
Surface pond & UV plus drying/refilling concentration & shallow well-mixed water body \\
Volcanic rock & wet/dry films and rock pores & porous surface rock/interstices \\
Hydrothermal vent & temperature, pH, redox, and flow gradients & porous chimney wall \\
\bottomrule
\end{longtable}

\section{First implementation target}

The first target is not a full ecological model. It is:
\begin{enumerate}
  \item hold a shallow pond state;
  \item update water depth in time;
  \item concentrate or dilute chemical species as water depth changes;
  \item add atmospheric source terms;
  \item apply UV photolysis and washout losses;
  \item expose the state to the chemical reaction network.
\end{enumerate}

\end{document}
